Apply \(\text{thm:adaptive-rootn-minimax}\).  Its once-for-all bank and positive dyadic certificate are computed from the certified-real primitive-constant names by \(\text{lem:finite-contour-bank}\), without access to the treatment-noise law.  On every index satisfying \(n\ge2\), \(\mathcal P_{\mathrm{NG},n}\neq\varnothing\), and \(\varepsilon_{1,n}\le\varepsilon_0\), that theorem gives constants \(0<c<C<\infty\), depending only on the fixed class constants, such that
\[
 \frac{c}{n}
 \le \mathfrak R_n^{\mathrm{NG}}
 \le \frac{C}{n}.
\]
The constants do not depend on \(n\).  Hence \(\mathfrak R_n^{\mathrm{NG}}\asymp n^{-1}\) along every sequence of indices in the statement.  In particular, the argument requires only \(\varepsilon_{1,n}\le\varepsilon_0\), not \(\varepsilon_{1,n}\to0\), so it includes every eligible fixed nonzero treatment-code radius.

The construction of \(\mathcal P_{\mathrm{NG},n}\) in \(\text{def:non-gaussian-class}\) contains the restriction \(\lVert\bar g_n-g_0\rVert_{L^1(P_X)}\le\varepsilon_{1,n}\) and contains no condition involving \(\widehat q_n\), \(\bar q_n\), or \(\varepsilon_{2,n}\).  Likewise, inspection of the explicit input list and algorithm in \(\text{def:adaptive-contour-estimator}\) shows that \(\widehat\theta_{\mathrm{spec},n}\) uses \(\bar g_n\) but none of those three outcome-code objects.  This proves the claimed absence of an outcome learner.

The two model classes are not being identified with one another.  Indeed, \(k=r+1\ge3\).  Every member of \(\mathcal P_{\mathrm G,n}^{\mathrm{JMS}}\) has Gaussian treatment noise, whose cumulants of orders at least three vanish, whereas every member of \(\mathcal P_{\mathrm{NG},n}\) satisfies \(\lvert\kappa_k(\eta)\rvert\ge\delta>0\).  Thus the classes are disjoint; whenever both are nonempty, neither contains the other.  No class inclusion is needed for either conclusion below.

Separately, \(\text{prop:bounded-outcome-gaussian-degeneracy}\) proves that every \(P\in\mathcal P_{\mathrm G,n}^{\mathrm{JMS}}\) satisfies \(\theta_0(P)=0\).  If this class is nonempty, the statistic identically equal to zero therefore has zero squared-error loss and zero absolute-error generalized-quantile loss for every member.  Since the losses in \(\text{def:minimax-risks}\) are nonnegative,
\[
 \mathfrak R_n^{\mathrm G}
 =\mathfrak M_{n,1-\gamma}^{\mathrm G}
 =0.
\]
The proposition derives this conclusion only from the simultaneous exact partially linear conditional-mean restriction, bounded observed outcome, and nondegenerate Gaussian treatment noise.  It invokes no positive lower-bound assertion from the source theorem.  Because the non-Gaussian fixed-separation rate and the Gaussian zero-target calculation concern separate model classes and supply no shrinking-separation comparison, they establish neither a Gaussian-versus-non-Gaussian minimax dichotomy nor a local-to-Gaussian frontier.  This proves all claims and scope qualifications.